\scanpage{32}
\begingroup\csname tagsleft@true\endcsname

% top left, two circled models
% slanted
\margin{\handwritten{\tikz[baseline=(c.base)]{\node[draw,ellipse,inner sep=1pt] (c) {$\phi^{(\ell)}$};}\ \tikz[baseline=(c.base)]{\node[draw,ellipse,inner sep=1pt] (c) {${}_{\ell^\nu}P$};}}}

% first formula
\[\deleted{\unsure{(}}\phi_{\handwritten{\xi}}^{(\ell)}\deleted{\unsure{)}}=(\phi_{\handwritten{\xi}}^{\ell^{\handwritten{\nu}}})_{\handwritten{\nu\geq}0}\,,\quad\text{où}\quad \phi_{\handwritten{\xi}}^{\ell^{\handwritten{\nu}}}:{}_{\ell^{\handwritten{\nu}}}P\;\mathsf{x}\;{}_{\ell^{\handwritten{\nu}}}Q\handwritten{\longrightarrow}G\,.\tag*{(2.2.4)}\]

% top, diagrams and systems
Les diagrammes\inserted{ commutatifs} déduits de (2.2.2) en y faisant $n=\ell^{\handwritten{\nu}}$, $n'=\ell^{\handwritten{\nu'}}$, pour $\handwritten{\nu'\geq\nu}$, montrent que $\phi_{\handwritten{\xi}}^{(\ell)}$ peut être interr̂été comme un accouplement de \deleted{pro-groupes} systèmes projectifs de groupes

% second formula
\[\phi_{\handwritten{\xi}}^{(\ell)}:\ T_\ell(P)\mathsf{x}\,T_\ell(Q)\handwritten{\longrightarrow}T_\ell(G)\,,\tag*{(2.2.5)}\]

% upper middle, introduction
où, pour un \deleted{\unsure{G}}groupe $M$ de $T$ et un nombre premier $\ell$, on désigne par

% projective system
\[T_\ell(M)=({}_{\ell^{\handwritten{\nu}}}M)_{\handwritten{\nu\geq}0}\]

% upper middle, name
le système p\deleted{\unsure{(}}rojectif des $\,{}_{\ell^{\handwritten{\nu}}}M$. L'accouplement \deleted{\unsure{$\phi$|à}} (2.2.5) est encore appelé \uline{l'accouplement} $\ell$-\uline{adique associé à} $\handwritten{\xi}$ (ou à la biextension $E$ de classe $\handwritten{\xi}$). Pour $\handwritten{\xi}$ variable, on obtient un homomorphisme canonique

% third formula
\[\mathrm{Biext}^{\handwritten{1}}(P,Q;G)\handwritten{\longrightarrow}\mathrm{Hom}(T_\ell(P)\otimes T_\ell(Q),T_\ell(G))\,.\tag*{(2.2.6)}\]

% middle, functoriality
On notera que les homomorphismes (2.1.2) sont \uline{fonctoriels} en $P,Q,G$ \uuline{(pour $n$} fixé), d'où on déduit qu'il en est de même de (2.2.6).

% section 2.3
2.3.\qquad \deleted{On peut regarder} La compatibilité (2.2.2) permet de retrouver $\phi_{\handwritten{\xi}}^n$ connaissant $\phi_{\handwritten{\xi}}^{n'}$, pourvu que $\,{}_{n'}Q\handwritten{\xrightarrow{\mathrm m}}{}_nQ$ \handwritten{soit} un épimorphisme, -- \deleted{(NB \unsure{il}} condition satisfaite, comme il est bien connu, si $nQ=Q$. En effet, comme $\,{}_{n'}P\handwritten{\xrightarrow{\mathrm m}}{}_nP$ est aussi un épimorphisme, $\phi_{\handwritten{\xi}}^n$ se déduit \deleted{\unsure{de}}alors de $\phi_{\handwritten{\xi}}^{n\handwritten{'}}$ par passage au quotient. On peut cependant \inserted{\handwritten{donner}} une compatibilité légèrement plus précise que (2.2.2), qui permet d\deleted{e}'exprimer $\phi^n$ en termes de $\phi^{n'}$ sans aucune restriction sur $Q$\handwritten{,} par la commutativité de

% lower square
\diagram{\begin{tikzcd}[ampersand replacement=\&,column sep=small]
\&{}_{n'}P\arrow[d,"\handwritten{m}"]\&\mathsf x\&{}_nQ\arrow[d,"\handwritten{\mathrm{id}}"]\arrow[rr,"{\phi_{\handwritten{\xi}}^{n'}\handwritten{|}\,{}_{\handwritten{n'}}P\;\mathsf{x}\;{}_nQ}"]\&\&{}_nG\arrow[d,"\handwritten{\mathrm{id}}"]\\
\text{(2.3.1)}\&{}_{\deleted{\unsure{m}}n}P\&\mathsf x\&{}_nQ\arrow[rr,"{\phi_{\handwritten{\xi}}^n}"]\&\&{}_nG\,,
\end{tikzcd}}

% under square
en formules

% penultimate formula
\[\phi_{\handwritten{\xi}}^n(mx,y)=\phi_{\handwritten{\xi}}^{n'}(x,y)\qquad(x\in{}_{n'}P(S),\ y\in{}_nQ(S)).\tag*{(2.3.2)}\]

% last lines
Cette formule implique évidemment la commutativité de (2.2.2), qui s'écrit en effet

% last formula
\[\phi_{\handwritten{\xi}}^n(mx,my)=m\,\phi_{\handwritten{\xi}}^{n'}(x,y)\qquad(x\handwritten{\in}{}_{n'}P(S),\ y\handwritten{\in}{}_{n'}P(S)),\tag*{\handwritten{(2.3.3)}}\]
\endgroup
% archive sheet number: 31
